\section{Appendix}\label{sec:appendix}

To prove Proposition~\ref{prop:strict},
we need the following Lemmata~\ref{lem:two-gen-nilpotent} and \ref{lem:two-gen-solvable}.

\begin{lemma}\label{lem:two-gen-nilpotent}
  For \(n \geq 1\), there is a two-generated finite nilpotent group of class \(n\).
\end{lemma}
\begin{proof}
  For each \(n \geq 1\), the dihedral group \(D_{2^{n + 1}}\) of order \(2^{n + 1}\)
  is a two-generated finite nilpotent group of class \(n\) (see \eg, \cite[p.\ 191, Example]{Dummit-Foote-2004}).
\end{proof}

\begin{lemma}\label{lem:two-gen-solvable}
  For \(n \geq 1\), there is a two-generated finite solvable group of derived length \(n\).
\end{lemma}
To prove Lemma~\ref{lem:two-gen-solvable}, we use the following embedding theorem due to Neumann-Neumann,
and Lemmata~\ref{lem:solvable-length-n} and \ref{lem:solvable-reduce} below.

\begin{theorem}[{\cite[Theorem 5.6]{neumann-neumann-1959}}]\label{thm:embed-solvable}
  Every finite solvable group \(G\) of derived length \(n\) can be embedded in
  some two-generated finite solvable group \(H\) of derived length at most \(n + 2\).
\end{theorem}

%\begin{lemma}[{\cite[Lemma 2.1]{AM2011}}]\label{lem:length-wreath}
%  Let \(G\) be a solvable group of derived length \(n\) and \(k \geq 2\) be an integer.
%  Then the wreath product \(G \wr C_k\) is a solvable group of derived length \(n + 1\),
%  where \(C_k\) is the finite cyclic group of order \(k\).
%\end{lemma}

\begin{lemma}\label{lem:solvable-length-n}
  For each \(n \geq 1\), there exists a finite solvable group of derived length \(n\).
\end{lemma}
\begin{proof}
  %By Lemma~\ref{lem:length-wreath},
  %the \(n\)-fold iterated wreath product \((\integer / 2 \integer) \wr \dotsb \wr (\integer / 2 \integer)\)
  %is a finite solvable group of derived length \(n\).
  For example, the \(n\)-fold iterated wreath product \((\integer / 2 \integer) \wr \dotsb \wr (\integer / 2 \integer)\)
  is a finite solvable group of derived length \(n\)~\cite[9.22 Lemma]{rose-group}.
\end{proof}

\begin{lemma}\label{lem:solvable-reduce}
  If \(G\) is a solvable group of derived length \(n + 1\), then some factor group of \(G\) is solvable of derived length \(n\).
\end{lemma}
\begin{proof}
  Let \(G = G^{(0)} \rhd G^{(1)} \rhd \dotsb \rhd G^{(n)} \rhd G^{(n + 1)} = \{1_G\}\) be the derived series of \(G\),
  where \(G^{(i + 1)} \defeq [G^{(i)}, G^{(i)}]\).
  Since each \(G^{(i + 1)}\) is a characteristic subgroup of \(G^{(i)}\), the group \(G^{(n)}\) is also a characteristic subgroup of \(G\),
  hence in particular a normal subgroup of each \(G^{(i)}\).
  Let \(\bar{G} \defeq G / G^{(n)}\).
  Then \(G^{(0)} / G^{(n)} \rhd G^{(1)} / G^{(n)} \rhd \dotsb \rhd G^{(n)} / G^{(n)} = \{1\}\) is the derived series of \(\bar{G}\) of length \(n\),
  since we have an isomorphism \([G^{(i)} / G^{(n)}, G^{(i)} / G^{(n)}] \cong [G^{(i)}, G^{(i)}] / G^{(n)} = G^{(i + 1)} / G^{(n)}\).
\end{proof}

\begin{proof}[Proof of Lemma~\ref{lem:two-gen-solvable}]
  By Lemma~\ref{lem:solvable-length-n}, there exists a finite solvable group \(G\) of derived length \(n\).
  By Theorem~\ref{thm:embed-solvable},
  \(G\) can be embedded in some two-generated finite group \(H\) of derived length at most \(n + 2\).
  %By Lemma~\ref{lem:length-wreath}, the derived length of \(G \wr C_k\) is \(n + 1\).
  In particular, the derived length of \(H\) is one of \(n\), \(n + 1\), or \(n + 2\).
  Thus, by Lemma~\ref{lem:solvable-reduce},
  we obtain a two-generated finite solvable group of derived length \(n\) as a factor group of \(H\).
\end{proof}

Lemma~\ref{lem:WP-synt} below is the last piece of the proof of Proposition~\ref{prop:strict}.
%Let \(G\) be a group and \(\eta\colon A^* \to G\) be a surjective morphism.
%The \emph{word problem} \(\WP_\eta(G)\) of \(G\) (with respect to \(\eta\))
%is the language over \(A\) defined as \(\WP_\eta(G) \defeq \eta^{-1}(1_G)\).
\begin{lemma}[folklore, \eg, {\cite[3.4.9 Proposition]{Holt2017}}]\label{lem:WP-synt}
  Let \(G\) be a group and \(\eta\colon A^* \to G\) be a surjective morphism.
  Then the syntactic monoid of the language \(\eta^{-1}(1_G) \subseteq A^*\) is isomorphic to \(G\).
  %Then the syntactic monoid of \(\WP_\eta(G)\) is isomorphic to \(G\).
  %In particular, every \(A\)-generated group is the syntactic monoid of some language over \(A\).
\end{lemma}
The language \(\eta^{-1}(1_G)\) is known as the \emph{word problem} for the group \(G\) (over \(A\)).

\begin{proof}[Proof of Proposition~\ref{prop:strict}]
  By Lemma~\ref{lem:two-gen-nilpotent}, there exists a two-generated finite nilpotent group \(G\) of class \(n + 1\).
  Since \(G\) is finite and generated by two elements as a group, it is generated by the two elements as a monoid.
  Hence there exists a surjective morphism \(\eta\colon A^* \to G\).
  Since \(\eta^{-1}(1_G)\) is clearly recognised by \(G\), we have \(\eta^{-1}(1_G) \in \gnil_A^{\leq n + 1}\).
  Contrary, suppose \(\eta^{-1}(1_G) \in \gnil_A^{\leq n}\).
  Then \(\eta^{-1}(1_G)\) is recognised by some finite nilpotent group \(H\) of class \(n\),
  hence the syntactic monoid of \(\eta^{-1}(1_G)\), which is isomorphic to \(G\) by Lemma~\ref{lem:WP-synt}, divides \(H\).
  In general, if \(K\) is a nilpotent group of class \(c\),
  then every subgroup of \(K\) and every factor group of \(K\) is nilpotent of class at most \(c\).
  Hence \(G\) is nilpotent of class at most \(n\), contradiction.
  (Note that every submonoid of a group is a subgroup.)
  Thus \(\gnil_A^{\leq n} \subsetneq \gnil_A^{\leq n + 1}\).

  By Lemma~\ref{lem:two-gen-solvable}, there exists a two-generated finite solvable group \(G\) of derived length \(n + 1\).
  Then one can similarly show \(\eta^{-1}(1_G) \in \gsol_A^{\leq n + 1} \setminus \gsol_A^{\leq n}\) for a surjective morphism \(\eta\colon A^* \to G\).
\end{proof}
